% Ad Atticum 13.11 --- headnote
% ad-atticum-13-11, 23 June 45 BC, Arpinum

\textit{Cicero to Atticus, written from the Arpinate
estate on the ninth day before the Kalends of Quintilis
709 AUC --- 23 June 45 BC (the manuscript dateline:
\textit{Scr.\ in Arpinati ix K.\ Quint.\ a.\ 709 (45)}).
The first letter from Arpinum after the move flagged in
13.10: Cicero had set out from the Tusculanum on or
about the eleventh before the Kalends, the trip being
forced on him both by the need to settle the rents of
his small estates and by a wish not to keep Brutus, who
was waiting at the next-door Tusculan villa, in
perpetual attendance on him.}

\textit{The letter opens with a snap of Greek ---
\textit{[Greek: ou tauton eidos]}, ``not the same
look'' --- the kind of clipped, almost proverbial
opening Cicero often uses when he is registering a
small private disappointment. The further distance from
Atticus and from Rome has altered the texture of the
days he had imagined would be straightforward. Section 2
is a short list of news he wants forwarded: any move by
Servilia (Brutus' mother, on the marital question
running through these letters), any action by Brutus
himself, the date Brutus has fixed for their meeting,
and, if convenient, an interview with Piso. The
\textit{maturum} of ripeness is left unexplained --- it
belongs to the same political matter Cicero and Atticus
have been turning over by letter, and the Arpinum
isolation makes him want it pushed forward.}
